\subsection{Coordinated finite bar lifts through a contractible parameter}

\begin{theorem}[C28: \status{Proved}; absolute finite bar lifts for projective parameters]
\label{sa08:C28}
Let $k$ be a commutative ring, put $S=\Gamma_k(T)$, let
$w>0$, let $q\geq1$ be odd, and let
$P$ be a simplicial $k$-module whose normalized terms are
projective. Let $A_1,\ldots,A_t$ be nonnegative bigraded
strict simplicial PD differential $S$-algebras, and prescribe maps
\[
 x_j\colon P\longrightarrow A_j^{w,q}
 \qquad(1\leq j\leq t).
\]
Suppose each of these target rows is horizontally acyclic.
Fix a finite arity bound $M\geq2$. For every $2\leq m\leq M$
and every $j$, suppose the coefficient map
\[
 \underline{k}\longrightarrow A_j^{mw,*}
\]
given by the weight-$mw$ coefficient $T^{[mw]}$
is an isomorphism on vertical $H^0$ and that vertical cohomology
vanishes in degrees $1,\ldots,m(q+1)$.

Then each prescribed original odd-to-even disk map extends
to a partial strict simplicial PD-DGA map from the positive
bar closure $U_q(P)$ through arity $M$. The maps may be
constructed together by factoring through one contractible
simplicial module $L$ with free normalized terms:
\[
 P\longrightarrow L,\qquad
 U_q(P)\longrightarrow U_q(L)\longrightarrow A_j.
\]
The last maps are understood through the stated finite arity.
No horizontal acyclicity of the target rows in the higher
arity weights is required.

If the vertical hypotheses hold for every $m$ at one fixed
$w$, compatible induction gives maps of the full bar closures.
The construction fixes the original disk maps; it does not
preserve arbitrary previously prescribed higher-arity values.
\end{theorem}
\begin{proof}
We first give a chain-level factorization with all its choices.
Set
\[
 W=\prod_{j=1}^t A_j^{w,q},\qquad C=N(W).
\]
The finite product is acyclic. For each $n\geq0$, choose a
free module $F_n$ surjecting onto $Z_n(C)$. Since $C$ is exact,
choose a lift of every basis element to $C_{n+1}$. These
choices define a chain map
\[
 E=\bigoplus_{n\geq0}D_{n+1}(F_n)\longrightarrow C,
\]
where $D_{n+1}(F_n)$ is the identity disk in homological
degrees $n+1,n$. This map is surjective in each degree.
Indeed, for $c\in C_n$, first lift $dc\in Z_{n-1}(C)$
using the upper part of $D_n(F_{n-1})$; the remaining
difference is a cycle and lifts using the lower part of
$D_{n+1}(F_n)$. The degree-zero assertion uses only the
second step. The complex $E$ is contractible and termwise
free. Its kernel $K$ is acyclic.

The combined chain map $N(P)\to C$ lifts to $E$.
Here is a direct construction, avoiding any inference from
termwise projectivity alone. Suppose the lift has been
constructed through degree $n-1$. Choose a graded lift
$g'_n\colon N_n(P)\to E_n$. Its differential defect
\[
 dg'_n-g_{n-1}d\colon N_n(P)\longrightarrow K_{n-1}
\]
takes values in $Z_{n-1}(K)$. Since
$K_n\to Z_{n-1}(K)$ is surjective and $N_n(P)$ is projective,
lift the defect to a map $h_n\colon N_n(P)\to K_n$.
Then $g_n=g'_n-h_n$ is the required chain lift.
At degree zero one simply chooses a lift. This induction
constructs the entire map. Put $L=\operatorname{DK}(E)$.
We have obtained the asserted factorization of all the
original parameter maps through a single $L$.

Now let $U_q(L)$ be the bar closure with coefficient
generators in form degree $q+1$. Arity $m$ and bar length
$n$ have form degree $m(q+1)-n$. Let $V_m(L)$ be the
arity-$m$ part of the PD subalgebra generated by smaller
arities, and set
\[
 Q_m(L)=U_q(L)_m/V_m(L).
\]
For each form degree $i$, the functor
$L\mapsto Q_m^i(L)$ is reduced and preserves simplicial
homotopies. It is a functor of modules even though it is
not an additive functor: the simplicial homotopy identities
involve compositions, and the reduced functor preserves
the zero map. Therefore $Q_m^i(L)$ is simplicially
contractible.

Its values in every simplex are free over $k$. Indeed,
$L$ is degreewise free, and Statement~\ref{sa08:C24}
identifies the quotient with the module on the signed
Lyndon PD generators of arity $m$. Normalized terms
are projective as direct summands of these free modules.
A contractible nonnegative complex of projective modules
is a projective object of the abelian category of such
complexes: its cycles are direct summands, and the complex
is a direct sum of identity disks on these projective
cycles. By Dold--Kan, $Q_m^i(L)$ is therefore a projective
object of $s\operatorname{Mod}_k$.

Start with the arity-one map defined by the chosen
$L\to A_j^{w,q}$ and its vertical differential.
Suppose a partial map has been constructed through
arity $m-1$. The free-subalgebra argument of
Statement~\ref{sa08:C25} gives its forced map on
$V_m(L)$, including every product and divided power
whose output has arity $m$. The obstruction to extending
it across $U_q(L)_m$ is an extension class in
\[
 \operatorname{Ext}^1_{\operatorname{Ch}(s\operatorname{Mod}_k)}
       \bigl(Q_m(L),A_j^{mw,*}\bigr).
\]
Since every term of $Q_m(L)$ is projective in
$s\operatorname{Mod}_k$, such extensions split in
each form degree, and the obstruction is represented in
\[
 H^1\operatorname{Tot}\operatorname{Hom}_{s\operatorname{Mod}_k}
       \bigl(Q_m(L),A_j^{mw,*}\bigr).
\]

The bounded complex $Q_m(L)[-1]$ consists of projective
objects and is supported through form degree
$m(q+1)$. The stated finite vertical comparison therefore
identifies this group with
\[
 H^1\operatorname{Hom}_{s\operatorname{Mod}_k}
       \bigl(Q_m(L),\underline{k}\bigr).
\]
This last Hom complex is identically zero: each
$Q_m^i(L)$ is horizontally acyclic, and
\[
 \operatorname{Hom}_{s\operatorname{Mod}_k}
       (Y,\underline{k})
 =\operatorname{Hom}_k(H_0NY,k)=0
\]
for such a $Y$. Thus the obstruction vanishes and an
actual simplicial cochain extension exists. All algebra
and PD identities in arity $m$ were already imposed on
$V_m(L)$, so the extension is a partial strict PD-DGA
map. Repeat for $m=2,\ldots,M$ and for every target.
Pulling back along $U_q(P)\to U_q(L)$ gives the required
maps and preserves every original $x_j$.

The same induction is countable when the target bounds
hold in every arity at a single weight. Taking the union
of compatible partial maps gives the asserted full maps.
\end{proof}

\paragraph{Why this remains an absolute result.}
The factorization uses that each prescribed parameter map
lands in an acyclic target row. It coordinates all later
choices through $L$. A preassigned map of the whole bar
closure on a parameter boundary need not factor this way.
In particular, a horizontal class in the source that a
relative attachment is meant to kill cannot already factor
through a contractible source parameter. The theorem also
does not choose one weight meeting an unbounded sequence
of eventual vertical thresholds.
